\chapter{Introduction}\label{ch:auto-002}

\section{What Computer Architecture Is}\label{sec:auto-003}
Computer architecture is the discipline of designing computing systems that execute real workloads efficiently under practical constraints.
In this course context, architecture means deciding how a processor and memory system should be organized so that useful work is completed quickly, reliably, and within power and cost limits.
It is important to separate this from neighboring fields.
Architecture is not transistor fabrication, circuit layout, or process chemistry, even though those influence what architects can build.

Every architectural decision is a tradeoff among competing goals.
A design that minimizes latency for one task may reduce throughput for a server-style workload.
A design that aggressively improves performance may consume more energy or require expensive cooling.
A design that is elegant on paper may be too costly in area or too difficult to verify at scale.
For this reason, architecture is best treated as constrained optimization: choose structures and policies that best match the target application while respecting implementation limits.

\section{Core Methods: Parallelism, Speculation, and Prediction}\label{sec:auto-004}
Most performance techniques in modern processors can be traced to three closely linked ideas: \textit{parallelism}, \textit{speculation}, and \textit{prediction}.
\textit{Parallelism} increases performance by doing multiple useful operations at the same time.
\textit{Speculation} increases performance by doing work early, before the machine is fully certain that this work will be needed.
\textit{Prediction} estimates uncertain future behavior (for example branch direction or branch target) so speculative work is directed toward the most likely path.
Prediction is not a third independent speedup method; it is a tool that makes speculation more selective and efficient.

\subsection{Speculation Without Prediction}\label{subsec:auto-001}
Speculation does not require a learned predictor.
A processor can speculate by optimistic overlap: continue execution under the assumption that an operation will complete normally, and recover only if that assumption fails.
This is useful because waiting for every operation to retire before moving on would underutilize the pipeline.

\paragraph{Example: write buffer}\label{par:auto-006}
A write buffer is a standard example.
When the core executes a store, it can enqueue the write and continue with subsequent instructions before that store reaches the lower memory hierarchy.
If the store later encounters a fault condition, architectural recovery and exception handling restore correctness.
The processor did not predict branch direction or data values in this case; it simply advanced optimistically to avoid stalling.

\subsection{Speculation With Prediction}\label{subsec:auto-002}
When prediction is added, speculation becomes more selective and usually more efficient.
The processor estimates likely future behavior, executes along the estimated path, and commits the work only if the estimate is correct.

\paragraph{Example: branch prediction}\label{par:auto-007}
Branch prediction is the canonical case.
If the predictor guesses that a branch is taken, instruction fetch and decode proceed down the predicted target path immediately.
A correct guess keeps the front end busy and reduces control-hazard stalls.
An incorrect guess forces the machine to squash the speculative instructions and restart on the correct path, which introduces a misprediction penalty.
The basic tradeoff is straightforward: high prediction accuracy yields high performance; low accuracy can erase the benefit of speculation.

\section{Parallelism by Grain Size}\label{sec:auto-005}
Parallelism is best understood by \textit{grain size}, meaning the amount of work grouped into one concurrent unit.
At the smallest scale, bit-level parallelism appears in wide datapaths and arithmetic logic that operate on many bit positions at once.
A 64-bit adder, for example, processes far more state per cycle than an 8-bit design because the hardware width itself exposes more fine-grained parallel work.

At \textit{instruction-level parallelism (ILP)}, the machine overlaps independent instructions within a thread.
Superscalar issue, out-of-order scheduling, and pipelining all exploit ILP\@.
Consider two arithmetic instructions: \texttt{R1=R2+R3} and \texttt{R4=R5+R6}.
Because neither depends on the other, a dual-issue machine can dispatch them together.
In contrast, \texttt{R1=R2+R3} followed by \texttt{R4=R1+R6} introduces a data dependence that limits overlap.
This illustrates the central ILP limitation: dependencies, control hazards, and hardware resource conflicts bound achievable parallel speedup.

At \textit{thread-level parallelism (TLP)}, multiple threads execute concurrently on multithreaded or multicore hardware.
TLP can provide substantial throughput gains, but only if the workload contains enough independent work and synchronization overhead remains controlled.
Inter-program parallelism is the coarsest common level, where independent processes run simultaneously under an operating system scheduler.
This is the dominant form in desktop multitasking and cloud services, where many programs share the same physical platform.

\section{Flynn's Taxonomy}\label{sec:auto-006}
Flynn's taxonomy classifies parallel machines by counting instruction streams and data streams.
In \textit{SISD} (single instruction, single data), one instruction stream acts on one data stream, matching the traditional scalar uniprocessor model.
In \textit{SIMD} (single instruction, multiple data), one instruction stream is applied across many data elements in lockstep, which is highly effective for dense data-parallel computations such as vector math, image processing, and many GPU kernels.

In \textit{MIMD} (multiple instruction, multiple data), multiple processing elements execute different instruction streams on different data streams.
This is the dominant model for modern multicore CPUs, distributed servers, and cluster systems because it supports general-purpose, irregular workloads.
\textit{MISD} (multiple instruction, single data) exists mostly in specialized contexts, such as certain fault-tolerant designs, and is uncommon in mainstream systems.

\subsection{SIMD vs. MIMD}\label{subsec:auto-003}
SIMD and MIMD solve different problems.
SIMD is efficient when the same operation must be applied repeatedly to large collections of data.
Its lockstep execution naturally provides synchronization across lanes, which reduces control overhead.
MIMD is more flexible because each processor can follow its own control path, but this flexibility introduces synchronization and communication costs that must be managed explicitly by hardware and software.

Modern GPUs often use a hybrid model sometimes called \textit{SIMT} (Single Instruction, Multiple Threads).
SIMT provides a SIMD-like execution efficiency for the hardware while allowing a MIMD-like programming model where each thread appears to have its own scalar program counter.
The hardware manages the divergence and convergence of these threads, grouping those on the same path into SIMD-style execution units.

\subsection{Where pipelining fits}\label{subsec:auto-004}
Pipelining is orthogonal to Flynn's taxonomy.
It is a temporal form of parallelism in which different stages operate on different instructions concurrently.
Because pipeline structure is an internal execution organization, pipelining can appear inside SISD cores, SIMD engines, and MIMD processors alike.

\section{Vectors: Packed-Word SIMD vs. Pipelined Vector Units}\label{sec:auto-007}
Packed-word SIMD and classic vector pipelines both exploit data parallelism, but they do so with different execution structures.
Packed-word SIMD divides a fixed-width register into lanes (for example, four 32-bit lanes in 128 bits) and applies one instruction across all lanes in lockstep.
An operation such as \texttt{ADD\_PACKED} performs multiple elementwise additions per instruction, and lane boundaries isolate carries so one lane's overflow does not spill into neighboring lanes.
This style underlies CPU SIMD extensions such as SSE, AVX, and NEON\@.

Vector pipelines, in contrast, stream long vectors through deep arithmetic pipelines.
A floating-point multiply can be viewed conceptually as operand preparation, multiply, and result normalization/writeback stages.
If a pipeline has $k$ stages and processes $n$ elements, a serial implementation needs $nk$ cycles, while a pipelined implementation needs approximately $k+n-1$ cycles after fill effects are accounted for (Eq.~\eqref{eq:auto-001}):
\begin{equation}\label{eq:auto-001}
\text{serial cycles} = nk,\qquad
\text{pipeline cycles} \approx k+n-1.
\end{equation}
The result is substantial throughput gain for long vectors even though single-element latency is not eliminated.

\paragraph{Chaining}\label{par:auto-008}
Vector chaining improves this further.
If a second vector operation consumes results from a first operation, the hardware can forward early produced elements directly into the consumer pipeline instead of waiting for the full vector to complete.
Chaining reduces producer-consumer delay and improves sustained throughput in vectorized loops.

\section{Common Chapter 1 Pitfalls}\label{sec:auto-008}
Students often conflate speculation with prediction.
The correct view is that prediction is one way to make speculation more effective, not the definition of speculation itself.
Another recurring error is to assume that additional cores guarantee proportional speedup; real speedup is limited by dependencies, communication, and serial regions of code.
It is also common to treat SIMD and MIMD as interchangeable, even though they target different workload structures and have different programming implications.

A final source of confusion is the difference between packed-word SIMD and long-vector pipelining.
Both process multiple data elements, but packed SIMD exploits lane-level parallel arithmetic within instruction-width registers, while vector pipelines exploit staged streaming over long element sequences.
Related to this, pipelining primarily improves throughput rather than the latency of a single operation.

\section{Chapter 1 Summary}\label{sec:auto-009}
Chapter 1 introduces computer architecture as disciplined performance engineering under practical constraints.
The chapter frames nearly all speedup mechanisms as forms of parallelism or speculation, then organizes parallelism by grain size from bits through instructions, threads, and entire programs.
Flynn's taxonomy provides a clean language for distinguishing SISD, SIMD, MISD, and MIMD machines, especially the practical differences between SIMD efficiency and MIMD flexibility.
Finally, the chapter distinguishes packed-word SIMD from pipelined vector execution and shows why chaining is valuable in vector data paths.
Together, these ideas form the conceptual base used throughout later note volumes.
